% Hafta 1 — Bölümleme: güvenilir hesaplama tabanını küçült
% Derleme: py -3.12 tools/latex/derle.py docs/week-1/assets/h01-10-bolumleme.tex
\documentclass[tikz,border=8pt]{standalone}
\usepackage{cen429-cizim}
\begin{document}
\begin{tikzpicture}[node distance=24mm and 24mm]
  \node[baslik] (b) at (0,0) {Bölümleme: güvenilir hesaplama tabanını küçült};
  \node[altbaslik] (alt) at (b.south west) {aralarından yalnız tutamaç ve şifreli veri geçer};

  \node[uyarikutu=64mm, below=8mm of alt.south west, anchor=north west, xshift=2mm, text width=56mm] (genel)
    {\kb{GENEL BÖLGE}\\[1mm]
     büyük · hızlı değişir\\[2mm]
     Arayüz\\
     İş mantığı\\[2mm]
     hata olasılığı YÜKSEK};
  \node[iyikutu=64mm, right=24mm of genel, text width=56mm] (cekirdek)
    {\kb{KORUNAN ÇEKİRDEK}\\[1mm]
     küçük · seyrek değişir\\[2mm]
     Dar arayüz (tutamaçlar)\\
     Kripto ve anahtar yönetimi\\
     Bütünlük / RASP denetimleri\\[2mm]
     incelenebilir, kanıtlanabilir};

  \path (genel.south) -- (cekirdek.south) coordinate[midway] (ortasi);
  \node[dip, text width=150mm, below=10mm of ortasi, anchor=north]
    {Güvenlik kritik kod ne kadar küçükse, doğruluğunu kanıtlamak o kadar kolaydır.};
\end{tikzpicture}
\end{document}
